\subsection{Ikosaedrische $\Phi$-Kaskade und Ra-Eichung}\label{sec:eabc-zeugen-phi}

Im algebraischen Sektor ($a\neq 0$, $M^2/\mathrm{Norm}^2\in\mathbb{Z}[\Phi]$) kalibriert der Urzeuge $(0,-1,0,0)$ mit $\mathrm{Norm}=\Phi$ und algebraischem Kern $\alpha=25610371200$ die thermische Rayleigh-Eichung
\[
  \Lambda_{\mathrm{therm}} = \frac{Ra_{\mathrm{krit}}}{\alpha_{\mathrm{Ur}}},
  \qquad
  Ra = \Lambda_{\mathrm{therm}}\cdot\alpha,
\]
gegen die klassische Bénard-Schwelle $Ra_{\mathrm{krit}}\approx 1708$. Es folgt $\Lambda_{\mathrm{therm}}\approx 6{,}669173\cdot 10^{-8}$.

\begin{table}[ht]
  \centering
  \caption{Kanonische Zeugen des $\Phi$-Sektors ($a\neq 0$) und $Ra$-Eichung mit $\Lambda_{\mathrm{therm}}\approx 6{,}669173\cdot 10^{-8}$.}
  \label{tab:eabc-phi}
  \begin{tabular}{rrrrr}
    \toprule
    $(c_e,c_a,c_b,c_c)$ & $\mathrm{Norm}$ & $\alpha$ & $Ra$ & Gitter-Verhältnis \\
    \midrule
    $(0, -1, 0, 0)$ & 1.61803 & $25610371200$ & 1708.00 & $Ra_{\mathrm{krit}}$ (Urzeuge) \\
    $(-1, -1, 0, 0)$ & 1.90211 & $7683111360$ & 512.40 & $\tfrac{3}{10}\,Ra_{\mathrm{krit}}$ \\
    $(-2, -1, -1, 0)$ & 2.76008 & $2186251200$ & 145.80 & $\tfrac{3}{35}\,Ra_{\mathrm{krit}}$ \\
    $(0, -2, 0, 0)$ & 3.23607 & $6402592800$ & 427.00 & $\tfrac{1}{4}\,Ra_{\mathrm{krit}}$ \\
    $(-2, -2, 0, 0)$ & 3.80423 & $1920777840$ & 128.10 & $\tfrac{3}{40}\,Ra_{\mathrm{krit}}$ \\
    $(-3, -2, 0, 0)$ & 4.41272 & $1061893440$ & 70.82 & $\tfrac{1}{24}\,Ra_{\mathrm{krit}}$ \\
    \bottomrule
  \end{tabular}
\end{table}

Die absteigenden $Ra$-Stufen brechen in exakten rationalen Brüchen zur kritischen Schwelle auf; dies deutet im \#Energiedoku-Programm auf quantisierte Energiedissipation entlang der Symmetrieachsen des 120-Zellers hin (vgl.\ Abschnitt~\ref{sec:eabc-phaenomenologie}).
